\section{14.2 NVIDIA Isaac Sim과 GPU 가속의 원리}
\begin{frame}{Isaac Sim: GPU로 수천 개를 동시에}
  \begin{itemize}
    \item \kb{Isaac Sim} (NVIDIA): Omniverse 플랫폼 위의 로봇
      시뮬레이션 환경 --- 물리 계산 \textbf{자체}를 GPU에서
      병렬로 수행
    \item 같은 로봇 환경을 CPU에서 하나씩 순차로 도는 대신,
      \textbf{수천 개의 복제본}을 GPU 위에서 동시에
    \item 사실적인 렌더링(카메라 센서 시뮬레이션)까지 ---
      ``\textbf{픽셀을 보고 판단하는}'' 로봇 정책 학습에도 널리 사용
  \end{itemize}
  \vskip0.6em
  \begin{block}{단, 이번 학기 실습에서는 시연 영상으로만}
    최소 사양: \textbf{RTX 4080 (16GB VRAM)}급 이상, \textbf{RT 코어}
    필요 --- RT 코어가 없는 A100\,·\,H100류 데이터센터 GPU는
    \textbf{미지원}. 학교의 GPU가 바로 H100/H200 계열이라 실습에서
    직접 돌릴 수 없다.
  \end{block}
\end{frame}
